% The receive buffer, one cell per 32 byte cache line, in the wrapped case.
\begin{tikzpicture}[font=\sffamily\small]
\node[font=\sffamily\bfseries\small, anchor=west] at (0mm,52mm)
  {rx\_buf, 512 bytes, drawn as sixteen cache lines};

% ---- the cache line the invalidate actually works on
\draw[tspan] (32mm,42mm) -- (40mm,42mm);
\node[font=\sffamily\scriptsize, anchor=west] at (43mm,42mm)
  {one cell is one 32 byte cache line, which is the unit the invalidate
   discards};

% ---- the two spans a wrap produces
\draw[tspan] (0mm,34mm) -- (24mm,34mm);
\node[tlabel] at (12mm,34mm) {span 2: 96 bytes};
\draw[tspan] (96mm,34mm) -- (128mm,34mm);
\node[tlabel] at (112mm,34mm) {span 1: 128 bytes};

% ---- the strip
\foreach \i in {3,...,11}
  {\node[mmres, minimum width=8mm, text width=6mm, minimum height=8mm]
     at ({\i*8mm},20mm) {};}
\foreach \i in {0,1,2,12,13,14,15}
  {\node[mmram, minimum width=8mm, text width=6mm, minimum height=8mm]
     at ({\i*8mm},20mm) {};}
\foreach \i/\b in {0/0, 4/128, 8/256, 12/384}
  {\node[font=\tiny, anchor=south] at ({\i*8mm},28.5mm) {\b};}
\node[font=\tiny, anchor=south] at (128mm,28.5mm) {512};

% ---- the two positions
\draw[flow] (24mm,8mm) -- (24mm,19.5mm);
\node[font=\sffamily\scriptsize, align=center] at (24mm,5mm)
  {pos = 512 minus NDTR = 96\\{\tiny read out of the engine, once}};
\draw[flow] (96mm,8mm) -- (96mm,19.5mm);
\node[font=\sffamily\scriptsize, align=center] at (96mm,5mm)
  {old\_pos = 384\\{\tiny where the last check stopped}};

\node[font=\sffamily\scriptsize, anchor=north west, text width=128mm] at (0mm,-4mm)
  {The engine has wrapped since the last check, so pos is below old\_pos and the
   new data is in two pieces. The consumer is handed each piece separately and
   never receives a span that crosses the end of the array. Because the buffer
   starts on a cache line and is a whole number of lines long, an invalidate
   over either span can never discard a line that also holds something else.};
\end{tikzpicture}
